%

%%%

%%%   Самарский государственный технический университет

%%%   Кафедра прикладной математики и информатики

%%%

%%%   Преамбула для статьи в журнал "Вестник Самарского государственного

%%%   технического университета. Серия: «Физико-математические науки»"

%%%   Ориентирована под MikTEx 2.4 for Win32

%%%

%%%   Хотя сам журнал собирается под GNU/Linux на дистрибутиве TeXLive



\documentclass[11pt,a4paper,twoside]{article}



\usepackage[cp1251]{inputenc}

\usepackage[T2A]{fontenc}

\usepackage[english,russian]{babel}

\usepackage{samgtubib}

\usepackage[dvips]{graphicx}

\usepackage[multiple]{footmisc}



\usepackage{samgtu}

\renewcommand{\VestnikYear}{2009} % Год издания Вестника

\renewcommand{\VestnikNumber}{2\,(19)} % Номер Вестника

\renewcommand{\VestnikMonth}{Ноябрь} % Месяц выхода Вестника

\renewcommand{\VestnikData}{01.11.09} % Дата подписания Вестника в печать

\renewcommand{\VestnikUPL}{26,64} % Усл. печ. л.

\renewcommand{\VestnikUIL}{25,73} % Уч.-изд. л.

\renewcommand{\VestnikRegNumber}{479/09} % Регистрационный номер

\renewcommand{\VestnikOrder}{994} % Номер заказа







%%%%

%%%   Самарский государственный технический университет

%%%   Кафедра прикладной математики и информатики

%%%

%%%   Преамбула для статьи в журнал "Вестник Самарского государственного

%%%   технического университета. Серия: «Физико-математические науки»"

%%%   Ориентирована под MikTEx 2.4 for Win32

%%%

%%%   Хотя сам журнал собирается под GNU/Linux на дистрибутиве TeXLive



\documentclass[11pt,a4paper,twoside]{article}



\usepackage[cp1251]{inputenc}

\usepackage[T2A]{fontenc}

\usepackage[english,russian]{babel}

\usepackage{samgtubib}

\usepackage[dvips]{graphicx}

\usepackage[multiple]{footmisc}



\usepackage{samgtu}

\renewcommand{\VestnikYear}{2009} % Год издания Вестника

\renewcommand{\VestnikNumber}{2\,(19)} % Номер Вестника

\renewcommand{\VestnikMonth}{Ноябрь} % Месяц выхода Вестника

\renewcommand{\VestnikData}{01.11.09} % Дата подписания Вестника в печать

\renewcommand{\VestnikUPL}{26,64} % Усл. печ. л.

\renewcommand{\VestnikUIL}{25,73} % Уч.-изд. л.

\renewcommand{\VestnikRegNumber}{479/09} % Регистрационный номер

\renewcommand{\VestnikOrder}{994} % Номер заказа







\usepackage{qrcode}

%
%\newcommand{\totop}[1]{\raisebox{6pt}[1pt][2pt]{#1}} 
%\newcommand{\tottop}[1]{\raisebox{12pt}[1pt][2pt]{#1}} 
%\newcommand{\totttop}[1]{\raisebox{18pt}[1pt][2pt]{#1}} 
%\newcommand{\tottttop}[1]{\raisebox{24pt}[1pt][2pt]{#1}} 
%\newcommand{\totttttop}[1]{\raisebox{30pt}[1pt][2pt]{#1}} 
\newcommand{\tobot}[1]{\raisebox{-6pt}[1pt][2pt]{#1}} 
%\newcommand{\tobbot}[1]{\raisebox{-12pt}[1pt][2pt]{#1}} 
%\newcommand{\tobbbot}[1]{\raisebox{-18pt}[1pt][2pt]{#1}}

\graphicspath{{./00PIC/}}

%\samgtupubl % для генерации pdf, отдаваемого в типографию СамГТУ
\pdfcrop % обрезка pdf для размещения на math-net.ru
\draft % На корректуре
\filllastpage % Последняя страница не заполнена не полностью
%\briefcomm % Статья является кратким сообщением

\vsgtu{1502}

\FirstPage{1}
\LastPage{12}
%
%
\begin{document}


\renewcommand{\baselinestretch}{0.95}

\hfill \qrcode[hyperlink,height=.8in]{http://dx.doi.org/10.14498/vsgtu1502}
\vspace{-.8in}


%%%%%%%%%%%%%%%%%%%%%%%%%%%%%%%%%%%%%%%%%%%%%%%%%%%%%%%%%%%%%%%%%%%%%%%%%%%%%%%%%%%%
%Информация о статье и об авторах на русском и английском языках

% Номер УДК
\udc{517.968.74}%Обработка текста. Подготовка текстов : Языки программирования. Компьютерные языки
\msc{34K10, 34K99}% Коды MSC см. http://www.ams.org/msc/

\rutitle{Обыкновенное интегро-дифференциальное уравнение с~вырожденным ядром и~интегральным~условием}{Обыкновенное интегро-дифференциальное уравнение\dots}

\entitle{An ordinary integro-differential equation with~a~degenerate kernel and an integral condition}

\ruauthor{Т.~К.~Юлдашев}{Ю\,л\,д\,а\,ш\,е\,в~Т.~К.}

\ruaffil{\rukrasnoyarsk.}
\enaffil{\enkrasnoyarsk.}
% Если несколько, то так  \institute{ Организация 1 \and Организация 2  \and Организация 3}

%\email{E-mail}{tursun.k.yuldashev@gmail.com} %E-mail's авторов


\ruauthordetails{1}{\fullauthor{\rupid{27151}{Турсун Камалдинович Юлдашев}}\regalia{к.ф.-м.н., доц.; \mem{tursun.k.yuldashev@gmail.com}}\position{доцент} \dept{каф. высшей математики}}

\enauthordetails{1}{\fullauthor{\enpid{27151}{Tursun K. Yuldashev}} \regalia{Cand. Phys. \& Math. Sci.; \mem{tursun.k.yuldashev@gmail.com}} \position{Associate Professor} \dept{Dept. of Higher Mathematics}}


\rukeywords{интегро-дифференциальное уравнение, вырожденное ядро, отражение аргумента, интегральное условие, однозначная разрешимость}

\ruabstract{Рассмотрены вопросы об однозначной разрешимости нелокальной краевой задачи для одного нелинейного обыкновенного интегро"=дифференциального уравнения с вырожденным ядром и отражением аргумента. Развит метод вырожденного ядра для случая рассматриваемого интегро"=дифференциального уравнения первого порядка. С помощью обозначения интегро"=дифференциальное уравнение сведено к системе алгебраических уравнений со сложной правой частью. После некоторого преобразования получено нелинейное функционально"=интегральное уравнение, однозначная разрешимость которого доказана методом последовательных приближений в сочетании его с методом сжимающих отображений. Данная работа является дальнейшим развитием теории обыкновенных нелинейных интегро-дифференциальных уравнений с вырожденным ядром.}
%

%
\enauthor{T.~K.~Yuldashev}{Y\,u\,l\,d\,a\,s\,h\,e\,v~T.~K.}


%\enginstitute{M. F. Reshetnev Siberian State Aerospace University, 31, pr. “Krasnoyarski Rabochiy”, Krasnoyarsk, 660014, Russia}% Место работы каждого из авторов на англ. языке
%% Если несколько, то так  \enginstitute{ Организация 1 \and Организация 2  \and Организация 3}
%

%
\enabstract{It is considered the questions of one value solvability of the nonlocal boundary value problem for a nonlinear ordinary integro-differential equation with degenerate kernel and reflective argument. The method of degenerate kernel is developed to the case of  considering ordinary integro-differential equation of the first order.  After denoting the integro-differential equation is reduced to a system of algebraic equations with complex right-hand side. After some transformation there is obtained the nonlinear functional-integral equation, one value solvability of which is proved by the aid of method of successive approximations combined it with the method of compressing mapping. This paper advances the theory of nonlinear integro-differential equations with a degenerate kernel.}

\enkeywords{integro-differential equation, degenerate kernel, reflective argument, integral form condition, one valued solvability}

\date{23/VII/2016} % Дата поступления статьи в редакцию.
\revisiondate{15/10/2016} % В окончательном виде
\accepteddate{09/XII/2016}

%%%%%%%%%%%%%%%%%%%%%%%%%%%%%%%%%%%%%%%%%%%%%%%%%%%%%%%%%%%%%%%%%%%%%%%%%%%%%%%%%%%%


\makerutitle

\Section{Постановка задачи}
Математическое моделирование многих процессов, происходящих в~реальном мире, часто приводит к~изучению начальных и~граничных задач для обыкновенных дифференциальных уравнений и~дифференциальных уравнений в~частных производных. Такие задачи составляют основу математической физики. 
Большой интерес с~точки зрения физических приложений представляют  интегро"=дифференциальные уравнения. 
Изучению обыкновенных интегро"=дифференциальных уравнений посвящено большое количество работ (см., например, 
[\citen{yul:11}--\citen{yul:100}]).
В~случаях, когда граница области протекания физического процесса недоступна для измерений, в~качестве дополнительной информации, достаточной для однозначной разрешимости задачи, могут служить нелокальные условия в~интегральной форме (см. [\citen{yul:111}--\citen{yul:1155}]). 

В~настоящей работе предлагается методика изучения нелокальной краевой задачи для обыкновенного нелинейного интегро"=дифференциального уравнения первого порядка с~вырожденным ядром. 
Методы исследования интегро"=дифференциальных уравнений в~частных производных с~вырожденным ядром развивались в~работах [\citen{yul:1166}--\citen{yul:2200}] 

Итак,  на отрезке
$ [-T ; T] $ рассмотрим интегро"=дифференциальное уравнение вида
\begin{equation}
u' (t)+\mu  \int _{-T}^T K  (t , s)  u (-s) \, d s=
\eta (t)  \int _{-T}^T u(\theta) \, d \theta+
f \biggl( \int _{-T}^T  H (\theta)  u (\theta) \, d \theta \biggr) \label{tur1}
\end{equation}
с интегральным условием
\begin {equation}
\label{tur2}
u (0)+ \int _{-T}^T \Theta (t)  u(t) \, d t=\varphi ,
\end{equation}
где $ \eta (t) \in C [-T ; T]$, $f (\gamma) \in C (\mathbb R) $,  
$ \Theta (t)  \in C^{ 1 } [-T ; T]$, 
$ \varphi=\rm const $, 
$$0< \int _{-T}^T | H (t) | \, d t<\infty ,$$ 
$$K (t , s)=\sum \limits_{i=1}^{m} a_{ i} (t)  b_{ i} (s) ,\quad  a_{i} (t) ,  b_{ i} (s) \in C^{ 1 }[-T ; T],$$ $0<T<\infty$,  $\mu $ 
"--- действительный спектральный параметр. 
Здесь предполагается, что функции $a_{i} (t)$ и $b_{i} (s)$ являются линейно независимыми.

	%\medskip

\smallskip
	
\Section{Нелинейное функционально-интегральное уравнение}
В уравнении \eqref{tur1} в интегральном слагаемом сделаем замену переменной: 
$$s=-\tau,\quad d s=-d \tau.$$ 
Тогда уравнение \eqref{tur1} приобретает вид
\begin{equation}
 u' (t) -\mu   \int _{-T}^T \sum \limits_{i=1}^{m} a_{i} (t)  b_{i} (-\tau)  u (\tau)  \, d \tau=\eta (t)  \int _{-T}^T u (\theta) \, d \theta+f (\gamma) , 
 \label{tur3}
\end{equation}
где $$\gamma= \int _{-T}^T  H (\theta)  u (\theta) \, d \theta .$$

С помощью обозначения 
\begin {equation}\label{tur4}
c_{i}= \int _{-T}^T b_{i} (-\tau)  u (\tau) \, d \tau
\end{equation}
уравнение (\ref{tur3}) перепишется в следующем виде:
\begin {equation}
\label{tur5}
 u' (t) =\mu  \sum \limits_{i=1}^{m} a_{i} (t)   c_{i}+\eta (t)  \int _{-T}^T u (\theta) \, d \theta+f (\gamma) .
\end{equation}
 
 Пусть 
 $$ \int _{-T}^T \Theta (t) \, d t \ne -1.$$ 
 Тогда	интегрированием из \eqref{tur5} с учетом интегрального условия 	\eqref{tur2} получаем 
\begin{multline}
u (t) =\varphi  h
-h  \int _{-T}^T \Theta (t) \biggl[ \mu  \sum \limits_{i=1}^{m} q_{i} (t)  c_{i}+p (t)  \int _{-T}^T u (\theta) \, d \theta+t  f (\gamma) \biggr] d t+
\\
+\mu  \sum \limits_{i=1}^{m} q_{i} (t)  c_{i}+p (t)  \int _{-T}^T u (\theta) \, d \theta+t  f (\gamma) , \label{tur6}
\end{multline}
где 
$$
h =\biggl( 1+ \int _{-T}^T \Theta (t) \, d t \biggr)^{ -1},\quad 
p (t)= \int _0^t \eta (\tau) \, d \tau , \quad    
q_{i} (t)= \int _0^t a_{i} (\tau) \, d \tau , \quad  i=\overline{1 , m} .$$
 
Подстановка выражения \eqref{tur6} в \eqref{tur4} дает систему алгебраических уравнений (САУ)
\begin {equation}\label{tur7}
c_{i}+\mu  \sum \limits_{j=1}^{m} A_{ i j}  c_{ j} =B_{i}  , \quad  i=\overline{1 , m} ,
\end{equation}
где
$$
A_{i j}=- \int _{-T}^T b_{i} (-\tau)  q_{j} (\tau) \, d \tau+
h   \int _{-T}^T b_{i} (-\tau)  \int _{-T}^T \Theta (\xi)  q_{ j} (\xi) \, d \xi  d \tau 
$$
и
\begin{multline}
B_{i}=h   \int _{-T}^T b_{i} (-\tau)  \int _{-T}^T \Theta (\xi) \biggl[ p (\xi)  \int _{-T}^T u (\theta) \, d \theta+\xi f (\gamma) \biggr] \, d \xi  d \tau-
\\
- \int _{-T}^T b_{i} (-\tau) \biggl[ \varphi  h+
p (\tau)   \int _{-T}^T u (\theta) \, d \theta+s  f (\gamma) \biggr] \, d \tau . 
\label{tur8}
\end{multline}
	
САУ \eqref{tur7} однозначно разрешима при любых конечных $B_{i} $,  если выполняется следующее условие:
\begin {equation}\label{tur9}
\Delta (\mu)=\left| \begin{array}{cccc}
1+\mu  A_{11} & \mu  A_{12}& \ldots & \mu  A_{1 m}\\
\mu  A_{21} & 1+\mu  A_{22} &  \ldots & \mu  A_{2 m} \\
\vdots &  \vdots &  \ddots & \vdots \\
\mu  A_{m 1} & \mu  A_{m 2} &   \ldots  & 1+\mu  A_{m m}
\end{array} \right| \ne 0 .
\end{equation}

Определитель $\Delta (\mu)$ в \eqref{tur9} есть многочлен относительно $\mu$ степени не выше $m$. 
Уравнение $\Delta (\mu)=0$ имеет не более $m$ различных корней. 
Эти корни являются собственными числами ядра интегро"=дифференциального уравнения \eqref{tur1}.  
Другие значения $\mu$ называются регулярными, при которых условие \eqref{tur9} выполняется. 
Для регулярных значений $\mu$ САУ \eqref{tur7} имеет единственное решение при любой конечной ненулевой правой части. 
В настоящей работе для таких регулярных значений параметра $\mu$  устанавливается однозначная разрешимость поставленной нелокальной задачи \eqref{tur1}, \eqref{tur2}. 
Тогда решения САУ  \eqref{tur7} записываются в виде
\begin {equation}\label{tur10}
c_{i}=\frac{\Delta_{i}(\mu)}{\Delta (\mu )} , \quad i=\overline{1 , m} , 
\end{equation}
где 
$
\Delta_{i} (\mu)=
$
$$
=\left| \begin{array}{ccccccc}
1+\mu  A_{11} & \ldots & \mu  A_{1 (i-1)} & B_{1}  & \mu  A_{1 (i+1)} & \ldots & \mu  A_{1 m}\\
\mu  A_{21} & \ldots & \mu  A_{2 (i-1)}  & B_{2}  & \mu  A_{2 (i+1)} &  \ldots & \mu  A_{2 m} \\
\vdots &  \vdots & \vdots & \vdots & \vdots &  \ddots & \vdots \\
\mu  A_{m 1} & \ldots & \mu  A_{m (i-1)} & B_{ m}  & \mu  A_{m (i+1)} & \ldots  & 1+\mu  A_{m m}
\end{array} \right|. 
$$

Среди элементов определителей $\Delta_{i} (\mu)$ находятся величины $B_{i} $.  В их составе, в свою очередь,  находится неизвестная функция $u (t)$. 
В~самом деле, эта неизвестная функция находилась в~правой части САУ \eqref{tur7}. 
Чтобы вывести эту функцию из знака определителей выражение в \eqref{tur8} перепишем в следующем виде:
$$
B_{i}=B_{ 1 i}+B_{ 2 i}  \int _{-T}^T u (\theta) \, d \theta+B_{ 3 i}   f (\gamma) ,
$$
где 
$$
B_{ 1 i}=\varphi  h  \int _{-T}^T b_{i} (-\tau) \, d \tau , \quad
B_{ 2 i}=  \int _{-T}^T b_{i} (-\tau) \biggl[ -p (\tau)+h  \int _{-T}^T \Theta (\xi) p (\xi) \, d \xi \biggr] \, d \tau , 
$$
$$
B_{3 i}=  \int _{-T}^T b_{i} (-\tau) \biggl[ -\tau+h   \int _{-T}^T \xi  \Theta (\xi) \, d \xi \biggr] \, d \tau .
$$
	
В этом случае, согласно свойствам определителя имеем
 $$
\Delta_{i} (\mu)=\Delta_{ 1 i} (\mu)+\Delta_{ 2 i} (\mu)   \int _{-T}^T u (\theta) \, d \theta + \Delta_{ 3 i} (\mu)  f (\gamma) ,
$$
где 
$
\Delta_{k i}(\mu)=
$
$$
=\left| \begin{array}{ccccccc}
1+\mu  A_{11} & \ldots & \mu  A_{1 (i-1)} & B_{k 1}  & \mu  A_{1 (i+1)} & \ldots & \mu  A_{1 m}\\
\mu  A_{21} & \ldots & \mu  A_{2 (i-1)}  & B_{k 2}  & \mu  A_{2 (i+1)} &  \ldots & \mu  A_{2 m} \\
\vdots &  \vdots & \vdots & \vdots & \vdots &  \ddots & \vdots \\
\mu  A_{m 1} & \ldots & \mu   A_{m (i-1)} & B_{k m}  & \mu  A_{m (i+1)} & \ldots  & 1+\mu  A_{m m}
\end{array} \right| ,$$
$
k=1 , 2, 3.$
Тогда формула \eqref{tur10} принимает вид
\begin {equation}\label{tur11}
c_{i}=\frac{\Delta_{ 1 i}(\mu)}{\Delta (\mu )}+
\frac{\Delta_{2 i}(\mu)}{\Delta (\mu )}   \int _{-T}^T u (\theta) \, d \theta
+\frac{\Delta_{ 3 i}(\mu)}{\Delta (\mu )}   f (\gamma) , \quad  i=\overline{1 , m} . 
\end{equation}

Подставляя \eqref{tur11} в \eqref{tur6}, получаем следующее нелинейное функционально"=интегральное уравнение (НФИУ):
\begin{equation}
u (t) =\Im_{1} (t;  u) \equiv Q (t)+ G (t)  \int _{-T}^T u (\theta) \, d \theta+ \Phi (t)  f \biggl( \int _{-T}^T  H (\theta)  u (\theta) d \theta \biggr) ,  \label{tur12}
\end{equation}
где
$$
Q (t)=h   \varphi-
\mu  h  \int _{-T}^T \Theta (t) \sum \limits_{i=1}^{m} q_{i} (t)  \frac{\Delta_{ 1 i}(\mu)}{\Delta (\mu )} \, d t+\mu  \sum \limits_{i=1}^{m} q_{i} (t)  \frac{\Delta_{1 i}(\mu)}{\Delta (\mu )} ,
$$
\begin{multline}
G (t)= p (t)-h  \int _{-T}^T \Theta (t)  p (t) \, d t-
\\
-\mu  h  \int _{-T}^T \Theta (t) \sum \limits_{i=1}^{m} q_{i} (t)  \frac{\Delta_{2 i}(\mu)}{\Delta (\mu )} \, d t+\mu \sum \limits_{i=1}^{m} q_{i} (t)  \frac{\Delta_{ 2 i}(\mu)}{\Delta (\mu )} ,
\end{multline}
\begin{multline}
\Phi (t)=t-h  \int _{-T}^T \Theta (t)  t \, d t-
\\
-\mu  h  \int _{-T}^T \Theta (t)  \sum \limits_{i=1}^{m} q_{i} (t)  \frac{\Delta_{ 3 i}(\mu)}{\Delta (\mu )} \, d t+\mu  \sum \limits_{i=1}^{m} q_{i} (t)  \frac{\Delta_{3 i}(\mu)}{\Delta (\mu )} .
\end{multline}

\smallskip

%\medskip

\Section{Однозначная разрешимость нелинейного функционально"=интегрального уравнения}
Рассмотрим множество функций 
$$ \bigl\{ u (t) \mid u (t)  \in C [-T ;  T] \bigr\}.$$
С~введением нормы 
$$
\bigl\| u (t) \bigr\|=\max_{t \in [-T ;  T]}  \bigl| u (t) \bigr|
$$
оно становится банаховым пространством.

	%\medskip

\smallskip

\hypertarget{yuld:th1}{}

\begin{theorem}[1]\it Пусть выполняется условие \eqref{tur9} и 
\begin{itemize}
\item[\rm 1)] $\beta_{ 1}= \| Q (t) \| <\infty;$
$\beta_{ 2}= \| G (t) \| <\infty;$
$\beta_{ 3}= \| \Phi (t) \| <\infty;$

\item[\rm 2)] $M= \| f (\gamma) \| <\infty;$
$| f (\gamma_1)-f (\gamma_2) | \le L | \gamma_1-\gamma_2 | ,$ $ 0<M ,$ $ L=\rm const ;$

\item[\rm 3)] $\rho=2  \beta_{ 2}  T+L  \delta_{ 2}  \beta_{ 3} <1;$ 
$\displaystyle \delta_{ 2}=  \int _{-T}^T  | H (t) | \, d t <\infty .$
\end{itemize}
Тогда НФИУ \eqref{tur12} имеет единственное решение на отрезке  $[-T ;  T].$ Это решение может быть найдено из следующего итерационного процесса{\/\rm:}
\begin {equation}
\label{tur13}
\begin{cases} 
u^{ 0} (t) =Q (t) , & \\
u^{ j+1} (t) =\Im_{ 1} (t;  u^{ j}), & j=0 , 1 , 2 , \ldots , \: t \in [-T ; \, T]. 
\end{cases}
\end{equation} 
\end{theorem}

\smallskip
	
\proof Рассмотрим шар $S ( u^{ 0} ; r_1 )$ с~радиусом 
$$r_1=\beta_1+\dfrac{\beta_3  M}{1-2  \beta_2  T}.$$ 
Для нулевого приближения, в силу первого условия теоремы~\hyperlink{yuld:th1}{1}, из  \eqref{tur13} имеем 
\begin {equation}
\label{tur14}
\| u^{ 0} (t) \| \le \beta_{1} .	
\end{equation}

В силу условий теоремы~\hyperlink{yuld:th1}{1} и неравенства \eqref{tur14} из \eqref{tur13} для первой разности получаем оценку 
\begin{equation}
\| u^{ 1} (t)-u^{ 0} (t) \| \le 2  
\| Q (t) \| \beta_2  T+
\beta_{ 3}  \| f (\gamma) \|=2  \beta_1  \beta_2  T+\beta_3  M .
\label{tur15}
\end{equation}

В силу условий теоремы~\hyperlink{yuld:th1}{1} и оценки \eqref{tur15} из \eqref{tur13} для разности $u^{ 2} (t)-u^{ 0} (t)$ получим
\begin{multline}
\| u^{ 2} (t)-u^{ 0} (t) \| \le
\\
\le 2  (2  \beta_1  \beta_2  T+\beta_3  M)  \beta_2  T+\beta_3  M= 
\beta_1 (2  \beta_2  T)^2+(2  \beta_2  T+1)  \beta_3  M .
\label{tur16}
\end{multline}

Далее, из \eqref{tur13} с учетом \eqref{tur16} имеем
\begin{multline}
\| u^{ 3} (t)-u^{ 0} (t) \| \le (2  \beta_1  \beta_2  T+\beta_3  M )  ( 2  \beta_2  T)^2+(2  \beta_2  T+1)  \beta_3  M=
\\
=\beta_1  (2  \beta_2  T)^3+\big( (2  \beta_2  T)^2+ 2  \beta_2  T+1 \big)  \beta_3  M .
\label{tur17}
\end{multline}

Продолжая этот процесс, аналогично \eqref{tur17} получаем
\begin{multline}
\| u^{ j} (t)-u^{ 0} (t) \| \le \beta_1  (2  \beta_2  T)^j
+\bigl( (2  \beta_2  T)^{j-1}+(2  \beta_2  T)^{j-2}+ \ldots + \\ +(2 \beta_2  T)^2+2  \beta_2  T+1 \bigr)  \beta_3  M.
\label{tur18}
\end{multline}

Из последнего условия теоремы~\hyperlink{yuld:th1}{1} следует, что $2 \beta_2  T<1$. 
Поэтому из \eqref{tur18}, переходя к пределу при $j \to \infty$, имеем
\begin{multline*}
\lim_{j \to \infty} \| u^{ j} (t)-u^{ 0} (t) \| \le  \\
\le
\lim_{j \to \infty} \Bigl[ \beta_1 ( 2  \beta_2  T)^j
+\bigl( (2  \beta_2  T)^{j-1}+(2  \beta_2  T)^{j-2}+ \ldots +\\
 +(2  \beta_2  T)^2+2  \beta_2  T+1 \bigr)  
\beta_3  M \Bigr],
\end{multline*}
откуда получаем оценку
\begin{equation}
\label{tur19}
\| u^{ \infty} (t)-u^{ 0} (t) \|<\beta_1+\frac{\beta_3  M}{1-2  \beta_2  T}=r_1 .
\end{equation}

Из \eqref{tur19} следует, что оператор в правой части \eqref{tur12} отображает шар $S ( u^{ 0} ; r_1 )$ в~себя.
Теперь для произвольной разности $u^{ j+1} (t)-u^{ j} (t)$ получим следующую оценку:
\begin{multline}
\| u^{ j+1} (t)-u^{ j} (t) \|
\le \left[2  \beta_{ 2}  T+L \delta_{ 2}  \beta_{ 3}  \right]  
\| u^{ j} (t)-u^{ j-1} (t) \|= 
\\
 =\rho  \| u^{ j} (t)-u^{ j-1} (t) \|< \| u^{ j} (t)-u^{ j-1} (t) \| .
 \label{tur20}
\end{multline}

В силу последнего условия теоремы~\hyperlink{yuld:th1}{1} из оценки \eqref{tur20} следует, что оператор в правой части \eqref{tur12} является сжимающим. 
Из оценок \eqref{tur14}, \eqref{tur19} и \eqref{tur20} мы заключаем, что для оператора \eqref{tur12} существует единственная неподвижная точка (см., напр.~\cite[стр. 389--401]{yul:2211}). 
Следовательно, на отрезке  $[-T ;  T]$ НФИУ \eqref{tur12} имеет единственное решение $ u (t)$. 
Кроме того, для скорости сходимости справедлива оценка
$$
\| u^{j+1} (t)-u (t) \| 
\le \frac{\rho^{j+1}}{1-\rho}  \bigl(2  \beta_1  \beta_2  T+\beta_3  M \bigr) .
$$
Дифференцируя НФИУ \eqref{tur12}, имеем
\begin{equation}
u' (t) =\Im_{2} (t; u) \equiv Q' (t)+ G' (t)  
\int _{-T}^T u (\theta) \, d \theta + 
\Phi' (t)  f \biggl( \int _{-T}^T H (\theta)  u (\theta) \, d \theta \biggr) ,  \label{tur21}
\end{equation}
где $Q' (t) \in C^{ 1 } [-T ;  T]$, 
$G' (t) \in C^{ 1 } [-T ;  T]$, 
$\Phi' (t) \in C^{ 1} [-T ;  T]$.~$~\square$

\smallskip


\hypertarget{yuld:th2}{}
  
\begin{theorem}[2]\it  Пусть выполняются условия теоремы {\rm \hyperlink{yuld:th1}{1}} и 
\begin{itemize}
\item[\rm 1)]
$N_{ 0}= \| Q' (t) \|<\infty;$
\item[\rm 2)]
$N_{ 1}= \| G' (t) \| \le \beta_{2};$
\item[\rm 3)]
$N_{ 2}= \| \Phi' (t) \| \le \beta_{ 3}.$
\end{itemize}
Тогда $ u' (t) \in C [-T ;  T].$
\end{theorem} 

\smallskip
	
\proof Рассмотрим шар 
$S \bigl( ({d}/{d t}) u^{ 0} ; r_2 \bigr)$ с радиусом 
$$
r_2=\max \Bigl\{ N_{ 0} ;  N_{ 1}  T  \Bigl(\beta_1+\dfrac{1}{1-2  \beta_2  T} \Bigr) +N_{ 2}  M \Bigr\}<\infty . 
$$ 

Для оператора \eqref{tur21} рассмотрим следующий итерационный процесс:
\begin {equation}\label{tur22}
\begin{cases} 
 \dfrac{d}{d t} u^{ 0} (t) =Q' (t) , & \\
 \dfrac{d}{d t} u^{ j+1} (t) =\Im_{ 2} (t; u^{ j}), & j=0 , 1 , 2 , \ldots . 
\end{cases}
\end{equation} 

Покажем, что последовательные приближения \eqref{tur22} не выходят из шара $S \bigl( ({d}/{d t}) u^{ 0} ; r_2 \bigr)$. 
В~силу первого условия теоремы~\hyperlink{yuld:th2}{2} из \eqref{tur22} для нулевого приближения имеем 
\begin {equation}\label{tur23}
\Bigl\|  \dfrac{d}{d t} u^{ 0} (t) \Bigr\| \le N_{0} .	
\end{equation}

Для первой разности в силу условий теоремы~\hyperlink{yuld:th2}{2} и неравенства \eqref{tur14} 
из \eqref{tur22} получаем оценку 
$$
\Bigl\|  \frac{d}{d t} u^{1} (t)- \frac{d}{d t} u^{0} (t) \Bigr\| \le 2  N_1  \beta_1  T+N_2  M .
$$

Для разности $ ({d}/{d t}) u^{2} (t)- ({d}/{d t}) u^{\, 0} (t)$ в силу условий теоремы~\hyperlink{yuld:th2}{2} и~оценки \eqref{tur15} из \eqref{tur22}  получим 
$$
\Bigl\|  \frac{d}{d t} u^{2} (t)- \frac{d}{d t} u^{0} (t) \Bigr\|
\le 2  (2  \beta_1  \beta_2  T+\beta_3  M)  N_1  T+N_2  M .
$$

Далее, из \eqref{tur22} с учетом \eqref{tur16} имеем
\begin{equation}
\Bigl\|  \frac{d}{d t} u^{3} (t)- \frac{d}{d t} u^{0} (t) \Bigr\|
\le N_1  T  \bigl( \beta_1  (2  \beta_2 T)^2+(2  \beta_2  T+1)  \beta_3  M \bigr)+N_2  M.
\label{tur24}
\end{equation}

Продолжая этот процесс, аналогично \eqref{tur24} получаем
\begin{multline}
\Bigl\|  \frac{d}{d t} u^{j} (t)- \frac{d}{d t} u^{0} (t) \Bigr\| \le \\ 
\le N_1  T  \Bigl[ \beta_1 (2  \beta_2  T)^j+
\bigl( (2   \beta_2   T)^{j-1}+(2   \beta_2  T)^{j-2}+ \ldots + \\ 
+(2  \beta_2  T)^2+2  \beta_2  T+1 \bigr)  \beta_3   M \Bigr]+N_2  M .
\label{tur25}
\end{multline}

Из последнего условия теоремы~\hyperlink{yuld:th1}{1} следует, что $2  \beta_2  T<1$. 
Поэтому из \eqref{tur25} при переходе к пределу при $j \to \infty$  имеем
\begin{equation}\label{tur26}
\Bigl\|  \frac{d}{d t} u^{\infty} (t)- \frac{d}{d t} u^{ 0} (t) \Bigr\|<  N_{1} T \Bigl(\beta_1+\frac{1}{1-2  \beta_2  T} \Bigr) +N_{ 2}  M \le r_2 .
\end{equation}

Из \eqref{tur23} и \eqref{tur26} заключаем, что оператор в правой части \eqref{tur21} отображает шар 
$S \bigl(({d}/{d t})  u^{ 0} ; r_2 \bigr)$ в~себя. 
Аналогично \eqref{tur20} для произвольной разности $({d}/{d t}) u^{ j+1} (t)- ({d}/{d t}) u^{ j} (t)$  справедлива следующая оценка: 
$$\Bigl\|  \frac{d}{d t} u^{ j+1} (t)- \frac{d}{d t} u^{ j} (t) \Bigr\|< \|  u^{ j} (t)- u^{  j-1} (t) \|.$$  
Кроме этого, для скорости сходимости также справедлива следующая оценка:
$$
\Bigl\| 
\frac{d}{d t} u^{ j+1} (t)- \frac{d}{d t} u (t) \Bigr\|  \le \frac{\rho^{ j+1}}{1-\rho}  (2 \beta_1  N_1  T+N_2  M ) .
$$
Отсюда следует, что $ u' (t) \in C [-T ;  T]$.~$~\square$
  
  

\thanks{{\bf Декларация о финансовых и других взаимоотношениях.}
Исследование не имело спонсорской поддержки.
Автор несет полную ответственность за предоставление окончательной версии рукописи в печать.
Окончательная версия рукописи была одобрена автором.
Автор не получал гонорар за статью.
}


\thanks{{\bf ORCID}

{\bf Турсун Камалдинович Юлдашев:} \url{http://orcid.org/0000-0002-9346-5362}

}


  
\RusRef

 \begin{thebibliography}{99}
 

\RBibitem{yul:11}
\by Банг~Н.~Д., Чистяков~В.~Ф., Чистякова~Е.~В.
\paper О некоторых свойствах вырожденных систем линейных интегро-дифференциальных уравнений.~I
\jour Изв. Иркутского гос. ун-та. Сер. Математика
\yr 2015
\vol 11
\pages 13--27
\mathnet{http://mi.mathnet.ru/iigum214}

\RBibitem{yul:22}
\by Быков~Я.~В.
\book О некоторых задачах теории интегро-дифференциальных уравнений
\publaddr Фрунзе
\publ Кирг. гос. ун-т
\yr 1957
\totalpages 327

\RBibitem{yul:33}
\by Вайнберг~М.~М.
\paper Интегро-дифференциальные уравнения
\serial Итоги науки. Сер. Мат. анал. Теор. вероятн. Регулир. 1962
\yr 1964
\pages 5--37
\publ ВИНИТИ
\publaddr М.
\mathnet{http://mi.mathnet.ru/intv82}
\mathscinet{http://www.ams.org/mathscinet-getitem?mr=178320}
\zmath{https://zbmath.org/?q=an:0184.33702}

\RBibitem{yul:44}
\by Васильев~В.~В.
\paper К вопросу о решении задачи Коши для одного класса линейных интегро-дифференциальных уравнений
\jour Изв. вузов. Матем.
\yr 1961
\issue 4
\pages 8--24
\mathnet{http://mi.mathnet.ru/ivm1896}
\mathscinet{http://www.ams.org/mathscinet-getitem?mr=143002}
\zmath{https://zbmath.org/?q=an:0112.06502}

\RBibitem{yul:55}
\by Виграненко~Т.~И.
\paper Об одном классе линейных интегро-дифференциальных уравнений
\jour Зап. Ленинград. горного ин-та
\yr 1956
\issue 3
\vol 33
\pages 176--186

\RBibitem{yul:66}
\by Власов~В.~В., Перез Ортиз~Р.
\paper Спектральный анализ интегро-дифференциальных уравнений, возникающих в~теории вязкоупругости и~теплофизике
\jour Матем. заметки
\yr 2015
\vol 98
\issue 4
\pages 630--634
\mathnet{http://mi.mathnet.ru/mz10829}
\crossref{10.4213/mzm10829}
\mathscinet{http://www.ams.org/mathscinet-getitem?mr=3438518}
\elib{http://elibrary.ru/item.asp?id=24850178}
%\transl
%\jour Math. Notes
%\yr 2015
%\vol 98
%\issue 4
%\pages 689--693
%\crossref{10.1134/S0001434615090357}
%\isi{http://gateway.isiknowledge.com/gateway/Gateway.cgi?GWVersion=2&SrcApp=PARTNER_APP&SrcAuth=LinksAMR&DestLinkType=FullRecord&DestApp=ALL_WOS&KeyUT=000363520200035}
%\scopus{http://www.scopus.com/record/display.url?origin=inward&eid=2-s2.0-84944896642}

\RBibitem{yul:77}
\by Кривошеин~Л.~Е.
\paper Об одном методе решения некоторых линейных интегро-дифференциальных уравнений
\jour Изв. вузов. Матем.
\yr 1960
\issue 3
\pages 168--172
\mathnet{http://mi.mathnet.ru/ivm2217}
\mathscinet{http://www.ams.org/mathscinet-getitem?mr=131739}
\zmath{https://zbmath.org/?q=an:0103.07701}

\RBibitem{yul:88}
\by Ландо~Ю.~К.
\paper Краевая задача для линейных интегро-дифференциальных уравнений типа Вольтерра в случае распадающихся краевых условий
\jour Изв. вузов. Матем.
\yr 1961
\issue 3
\pages 56--65
\mathnet{http://mi.mathnet.ru/ivm1879}
\mathscinet{http://www.ams.org/mathscinet-getitem?mr=166574}
\zmath{https://zbmath.org/?q=an:0167.12001}

\RBibitem{yul:99}
\by Шишкин~Г.~А.
\paper Обоснование одного метода решения интегро-дифференциальных уравнений Фредгольма с~функциональным запаздыванием
\inbook Корректные краевые задачи для неклассических уравнений математической физики
\publaddr Новосибирск
\publ Ин-т математики СО АН СССР
\yr 1980
\pages 172--178

\RBibitem{yul:100}
\by Фалалеев~М.~В.
\paper Интегро-дифференциальные уравнения с фредгольмовым оператором при старшей производной в банаховых пространствах и их приложения
\jour Изв. Иркутского гос. ун-та. Сер. Математика
\yr 2012
\vol 5
\issue 2
\pages 90--102
\mathnet{http://mi.mathnet.ru/iigum70}

\RBibitem{yul:111}
\by Гордезиани~Д.~Г., Авалишвили~Г.~А.
\paper Решения нелокальных задач для одномерных колебаний среды
\jour Матем. моделирование
\yr 2000
\vol 12
\issue 1
\pages 94--103
\mathnet{http://mi.mathnet.ru/mm832}
\mathscinet{http://www.ams.org/mathscinet-getitem?mr=1773208}
\zmath{https://zbmath.org/?q=an:1027.74505}

\RBibitem{yul:1122}
\by Тихонов~И.~В.
\paper Теоремы единственности в~линейных нелокальных задачах для абстрактных дифференциальных уравнений
\jour Изв. РАН. Сер. матем.
\yr 2003
\vol 67
\issue 2
\pages 133--166
\mathnet{http://mi.mathnet.ru/izv429}
\crossref{10.4213/im429}
\mathscinet{http://www.ams.org/mathscinet-getitem?mr=1972995}
\zmath{https://zbmath.org/?q=an:1073.34071}
%\transl
%\jour Izv. Math.
%\yr 2003
%\vol 67
%\issue 2
%\pages 333--363
%\crossref{10.1070/IM2003v067n02ABEH000429}
%\isi{http://gateway.isiknowledge.com/gateway/Gateway.cgi?GWVersion=2&SrcApp=PARTNER_APP&SrcAuth=LinksAMR&DestLinkType=FullRecord&DestApp=ALL_WOS&KeyUT=000185541900005}
%\scopus{http://www.scopus.com/record/display.url?origin=inward&eid=2-s2.0-33748696820}

\RBibitem{yul:1133}
\by Пулькина~Л.~С.
\paper Нелокальная задача для гиперболического уравнения с~интегральными условиями I~рода с~ядрами, зависящими от времени
\jour Изв. вузов. Матем.
\yr 2012
\issue 10
\pages 32--44
\mathnet{http://mi.mathnet.ru/ivm8743}
\mathscinet{http://www.ams.org/mathscinet-getitem?mr=3137057}
%\transl
%\jour Russian Math. (Iz. VUZ)
%\yr 2012
%\vol 56
%\issue 10
%\pages 26--37
%\crossref{10.3103/S1066369X12100039}
%\scopus{http://www.scopus.com/record/display.url?origin=inward&eid=2-s2.0-84869422481}

\RBibitem{yul:1144}
\by Сабитова~Ю.~К.
\paper Краевая задача с~нелокальным интегральным условием для уравнений смешанного типа с~вырождением
на переходной линии
\jour Матем. заметки
\yr 2015
\vol 98
\issue 3
\pages 393--406
\mathnet{http://mi.mathnet.ru/mz9135}
\crossref{10.4213/mzm9135}
\mathscinet{http://www.ams.org/mathscinet-getitem?mr=3438496}
\elib{http://elibrary.ru/item.asp?id=24073750}
%\transl
%\jour Math. Notes
%\yr 2015
%\vol 98
%\issue 3
%\pages 454--465
%\crossref{10.1134/S0001434615090114}
%\isi{http://gateway.isiknowledge.com/gateway/Gateway.cgi?GWVersion=2&SrcApp=PARTNER_APP&SrcAuth=LinksAMR&DestLinkType=FullRecord&DestApp=ALL_WOS&KeyUT=000363520200011}
%\scopus{http://www.scopus.com/record/display.url?origin=inward&eid=2-s2.0-84944897163}

\RBibitem{yul:1155}
\by Тагиев~Р.~К., Габибов~В.~М.
\paper Об одной задаче оптимального управления для уравнения теплопроводности с интегральным граничным условием 
\jour Вестн. Сам. гос. техн. ун-та. Сер. Физ.-мат. науки
\yr 2016
\vol 20
\issue 1
\pages 54--64
\crossref{10.14498/vsgtu1463}

\RBibitem{yul:1166}
\by Юлдашев~Т.~К.
\paper Обратная задача для одного интегро-дифференциального уравнения Фредгольма в частных производных третьего порядка
\jour Вестн. Сам. гос. техн. ун-та. Сер. Физ.-мат. науки
\yr 2014
\issue 1(34)
\pages 56--65
\crossref{10.14498/vsgtu1299}

\RBibitem{yul:1177}
\by Юлдашев~Т.~К.
\paper Двойная обратная задача для интегро-дифференциального уравнения Фредгольма эллиптического типа
\jour Вестн. Сам. гос. техн. ун-та. Сер. Физ.-мат. науки
\yr 2014
\issue 2(35)
\pages 39--49
\crossref{10.14498/vsgtu1306}

\RBibitem{yul:1188}
\by Юлдашев~Т.~К.
\paper Об одном интегро-дифференциальном уравнении Фредгольма в~частных производных третьего порядка
\jour Изв. вузов. Матем.
\yr 2015
\issue 9
\pages 74--79
\mathnet{http://mi.mathnet.ru/ivm9038}
%\transl
%\jour Russian Math. (Iz. VUZ)
%\yr 2015
%\vol 59
%\issue 9
%\pages 62--66
%\crossref{10.3103/S1066369X15090091}
%\scopus{http://www.scopus.com/record/display.url?origin=inward&eid=2-s2.0-84943244664}

\RBibitem{yul:1199}
\by Юлдашев~Т.~К.
\paper Обратная задача для нелинейного интегро- дифференциального уравнения Фредгольма четвертого порядка с вырожденным ядром
\jour Вестн. Сам. гос. техн. ун-та. Сер. Физ.-мат. науки
\yr 2015
\vol 19
\issue 4
\pages 736--749
\crossref{10.14498/vsgtu1434}

\RBibitem{yul:2200}
\by Юлдашев~Т.~К.
\paper Обратная задача для интегро-дифференциального уравнения Фредгольма третьего порядка с~вырожденным ядром
\jour Владикавк. матем. журн.
\yr 2016
\vol 18
\issue 2
\pages 76--85
\mathnet{http://mi.mathnet.ru/vmj583}

\RBibitem{yul:2211}
\by Треногин~В.~А.
\book Функциональный анализ
\publaddr М.
\publ Наука
\yr 1980
\totalpages 495

\end{thebibliography}

\RuDate

\newpage


\makeentitle




\thanks{{\bf Declaration of Financial and Other Relationships.}
The research has not had any sponsorship.
The author is absolutely responsible for submitting the final manuscript in print.
The author has approved the final version of manuscript.
The author has not received any fee for the article.
}


\thanks{{\bf ORCID}

{\bf Tursun K. Yuldashev:} \url{http://orcid.org/0000-0002-9346-5362}
}

\newpage

\EngRef	

\begin{thebibliography}{99}
 

\Bibitem{yul:11:en}
\lang In Russian
\by Bang~N.~D., Chistyakov~V.~F., Chistyakova~E.~V.
\paper About some properties of degenerate systems of linear integro-differential equations.~I
\jour IIGU Ser. Matematika
\yr 2015
\vol 11
\pages 13--27

\Bibitem{yul:22:en}
\lang In Russian
\by Bykov~Ya.~V.
\book O nekotorykh zadachakh teorii integro-differentsial'nykh uravnenii \rm [On some problems in the theory of integrodifferential equations]
\publaddr Frunze
\publ Kirgiz. State Univ.
\yr 1957
\totalpages 327


\Bibitem{yul:33:en}
\lang In Russian
\by Vainberg~M.~M.
\paper Integro-differential equations
\serial Itogi Nauki. Ser. Mat. Anal. Teor. Ver. Regulir. 1962
\yr 1964
\pages 5--37
\publ VINITI
\publaddr Moscow


\Bibitem{yul:44:en}
\lang In Russian
\by Vasil'ev~V.~V.
\paper On the solution of the Cauchy problem for a class of linear integro-differential equations
\jour Izv. Vyssh. Uchebn. Zaved. Mat.
\yr 1961
\issue 4
\pages 8--24

\Bibitem{yul:55:en}
\lang In Russian
\by Vigranenko~T.~I.
\paper A boundary value problem for linear integro-differential equations
\jour Zap. Leningrad. gornogo in-ta
\yr 1956
\issue 3
\vol 33
\pages 176--186


\Bibitem{yul:66:en}
\paper Spectral analysis of integro-differential equations in viscoelasticity and thermal physics
\by Vlasov~V.~V.,  Perez Ortiz~R.
\jour Math. Notes
\yr 2015
\vol 98
\issue 4
\pages 689--693
\crossref{10.1134/S0001434615090357}
\isi{http://gateway.isiknowledge.com/gateway/Gateway.cgi?GWVersion=2&SrcApp=PARTNER_APP&SrcAuth=LinksAMR&DestLinkType=FullRecord&DestApp=ALL_WOS&KeyUT=000363520200035}
\scopus{http://www.scopus.com/record/display.url?origin=inward&eid=2-s2.0-84944896642}

\Bibitem{yul:77:en}
\lang In Russian
\by Krivoshein~L.~E.
\paper A~method of solving certain linear integro-differential equations
\jour Izv. Vyssh. Uchebn. Zaved. Mat.
\yr 1960
\issue 3
\pages 168--172

\Bibitem{yul:88:en}
\lang In Russian
\by Lando~Yu.~K.
\paper A boundary-value problem for linear integro-differential equations of Volterra type in the case of disjoint boundary conditions
\jour Izv. Vyssh. Uchebn. Zaved. Mat.
\yr 1961
\issue 3
\pages 56--65

\Bibitem{yul:99:en}
\lang In Russian
\by Shishkin~G.~A.
\paper Justification of a method for solving integro-differential Fredholm equations with functional delay
\inbook Korrektnye kraevye zadachi dlia neklassicheskikh uravnenii matematicheskoi fiziki \rm [Correct boundary value problems for nonclassical equations of mathematical physics]
\publaddr Novosibirsk
\publ Institute of Mathematics, Siberian Branch of the USSR Academy of Sciences
\yr 1980
\pages 172--178


\Bibitem{yul:100:en}
\lang In Russian
\by Falaleev~M.~V.
\paper Integro-differential equations with Fredholm operator by the derivative of the higest order in Banach spaces and it's applications
\jour IIGU Ser. Matematika
\yr 2012
\vol 5
\issue 2
\pages 90--102

\Bibitem{yul:111:en}
\lang In Russian
\by Gordeziani~D.~G., Avalishvili~G.~A.
\paper On the constructing of solutions of the nonlocal initial boundary value problems for one-dimensional medium oscillation equations
\jour Matem. Mod.
\yr 2000
\vol 12
\issue 1
\pages 94--103

\Bibitem{yul:1122:en}
\paper Uniqueness theorems for linear non-local problems for abstract differential equations
\by Tikhonov~I.~V.
\jour Izv. Math.
\yr 2003
\vol 67
\issue 2
\pages 333--363
\crossref{10.1070/IM2003v067n02ABEH000429}
\isi{http://gateway.isiknowledge.com/gateway/Gateway.cgi?GWVersion=2&SrcApp=PARTNER_APP&SrcAuth=LinksAMR&DestLinkType=FullRecord&DestApp=ALL_WOS&KeyUT=000185541900005}
\scopus{http://www.scopus.com/record/display.url?origin=inward&eid=2-s2.0-33748696820}

\Bibitem{yul:1133:en}
\paper A nonlocal problem for a hyperbolic equation with integral conditions of the 1st kind with time-dependent kernels
\by Pul'kina~L.~S.
\jour Russian Math. (Iz. VUZ)
\yr 2012
\vol 56
\issue 10
\pages 26--37
\crossref{10.3103/S1066369X12100039}
\scopus{http://www.scopus.com/record/display.url?origin=inward&eid=2-s2.0-84869422481}

\Bibitem{yul:1144:en}
\paper Boundary-value problem with nonlocal integral condition for mixed-type equations with degeneracy on the transition line
\by Sabitova~Yu.~K.
\jour Math. Notes
\yr 2015
\vol 98
\issue 3
\pages 454--465
\crossref{10.1134/S0001434615090114}
\isi{http://gateway.isiknowledge.com/gateway/Gateway.cgi?GWVersion=2&SrcApp=PARTNER_APP&SrcAuth=LinksAMR&DestLinkType=FullRecord&DestApp=ALL_WOS&KeyUT=000363520200011}
\scopus{http://www.scopus.com/record/display.url?origin=inward&eid=2-s2.0-84944897163}

\Bibitem{yul:1155:en}
\lang In Russian
\by Tagiyev~R.~K., Gabibov~V.~M.
\paper On optimal control problem for the heat equation with integral boundary condition 
\jour  Vestn. Samar. Gos. Tekhn. Univ. Ser. Fiz.-Mat. Nauki \rm [J. Samara State Tech. Univ., Ser. Phys. \& Math. Sci.]
\yr 2016
\vol 20
\issue 1
\pages 54--64
\crossref{10.14498/vsgtu1463}

\Bibitem{yul:1166:en}
\lang In Russian
\by Yuldashev~T.~K.
\paper Inverse Problem for a Fredholm Third Order Partial Integro-differential Equation
\jour  Vestn. Samar. Gos. Tekhn. Univ. Ser. Fiz.-Mat. Nauki \rm [J. Samara State Tech. Univ., Ser. Phys. \& Math. Sci.]
\yr 2014
\issue 1(34)
\pages 56--65
\crossref{10.14498/vsgtu1299}

\Bibitem{yul:1177:en}
\lang In Russian
\by Yuldashev~T.~K.
\paper A double inverse problem for Fredholm integro-differential equation of elliptic type
\jour  Vestn. Samar. Gos. Tekhn. Univ. Ser. Fiz.-Mat. Nauki \rm [J. Samara State Tech. Univ., Ser. Phys. \& Math. Sci.]
\yr 2014
\issue 2(35)
\pages 39--49
\crossref{10.14498/vsgtu1306}

\Bibitem{yul:1188:en}
\by Yuldashev~T.~K.
\paper A certain Fredholm partial integro-differential equation of the third order
\jour Russian Math. (Iz. VUZ)
\yr 2015
\vol 59
\issue 9
\pages 62--66
\crossref{10.3103/S1066369X15090091}
\scopus{http://www.scopus.com/record/display.url?origin=inward&eid=2-s2.0-84943244664}

\Bibitem{yul:1199:en}
\lang In Russian
\by Yuldashev~T.~K.
\paper An inverse problem for a nonlinear Fredholm   
integro-differential equation of fourth order with degenerate kernel
\jour  Vestn. Samar. Gos. Tekhn. Univ. Ser. Fiz.-Mat. Nauki \rm [J. Samara State Tech. Univ., Ser. Phys. \& Math. Sci.]
\yr 2015
\vol 19
\issue 4
\pages 736--749
\crossref{10.14498/vsgtu1434}

\Bibitem{yul:2200:en}
\lang In Russian
\by Yuldashev~T.~K.
\paper Inverse problem for a~third order Fredholm integro-differential equation with degenerate kernel
\jour Vladikavkaz. Mat. Zh.
\yr 2016
\vol 18
\issue 2
\pages 76--85

\Bibitem{yul:2211:en}
\lang In Russian
\by Trenogin~V.~A.
\book Funktsional'nyi analiz \rm [Functional analysis]
\publaddr Moscow
\publ Nauka
\yr 1980
\totalpages 495



\end{thebibliography}

\EngDate

\renewcommand{\baselinestretch}{0.9}

\end{document}
